\Askhsh{34332}{2}
\begin{ylh}{1}
    \item 3.4. 3ο Κριτήριο ισότητας τριγώνων
    \item 4.2. Τέμνουσα δύο ευθειών - Ευκλείδειο αίτημα
    \item 5.6. Εφαρμογές στα τρίγωνα.
\end{ylh}
% ===================================================================
%  Αρχείο: 34332.tex (Fragment)
%  Αυτόματη μετατροπή μέσω doc2latex.py
% ===================================================================

ΘΕΜΑ 2

Στο παρακάτω σχήμα δίνονται τα ισοσκελή τρίγωνα ΑΒΓ και ΕΓΔ με $ΑΒ = Α\Gamma = Ε\Gamma = Ε\Delta$, όπου Δ είναι το μέσο της ΑΓ και $Β\Gamma = \frac{ΑΒ}{2}$. Έστω Ζ το σημείο στο οποίο η προέκταση της ΕΔ προς το Δ τέμνει την ΑΒ.

Να αποδείξετε ότι:

\begin{alist}
\item τα τρίγωνα ΑΒΓ και ΓΔΕ είναι ίσα, \monades{12}

\item το σημείο Ζ είναι το μέσο της ΑΒ. \monades{13}


\begin{center}
\begin{tikzpicture}[scale=1]
% BC = AB/2. Let's make BC=2, AB=AC=4.
\tkzDefPoint(0,0){B}
\tkzDefPoint(2,0){C}
\tkzDefTriangle[two angles = 75 and 75](B,C) \tkzGetPoint{A} % wait, if AB=4, BC=2. Base is 2, legs 4.
% actually we can construct it exactly
\tkzInterCC[R](B,4)(C,4) \tkzGetPoints{A}{A2}

\tkzDefMidPoint(A,C) \tkzGetPoint{D}

% Triangle CDE is isosceles with EC=ED=AC=4.
% E is some point such that EC=ED=4.
% Wait: "ΕΓΔ με AB=AC=EΓ=EΔ". So EC=ED=4. D is midpoint, so CD=2.
% We have triangle CDE with base CD=2, legs EC=ED=4.
% Since AC=4 and CD=2, this is similar to ABC!
% Actually, it is "τα ισοσκελή τρίγωνα ΑΒΓ και ΕΓΔ".
% Where does E lie? Usually on the opposite side of AC or something. Let's place E.
\tkzInterCC[R](C,4)(D,4) \tkzGetPoints{E_alt}{E} % E to the right

% Extension of ED towards D intersects AB at Z.
\tkzInterLL(E,D)(A,B) \tkzGetPoint{Z}

\draw[l3] (A)--(B)--(C)--cycle;
\draw[l3] (E)--(C)--(D)--cycle;
\draw[l3] (E)--(Z);
\draw[l3] (A)--(B); % AZB line

\tkzDrawPoints(A,B,C,D,E,Z)
\tkzLabelPoint[above](A){$A$}
\tkzLabelPoint[left](B){$B$}
\tkzLabelPoint[below](C){$\Gamma$}
\tkzLabelPoint[right](D){$\Delta$}
\tkzLabelPoint[right](E){$E$}
\tkzLabelPoint[left](Z){$Z$}

\tkzMarkSegments[mark=s|](A,B A,C E,C E,D)
\tkzMarkSegments[mark=s||](A,D D,C B,C) % since AD=DC=2, BC=2
\end{tikzpicture}
\end{center}
\end{alist}
